% RESOURCE-DIRECTSUM REPORT::RESOURCE-PROMISE-ENTROPY12::2026-10-18_2330_KST
\documentclass{article}
\usepackage{amsmath,amsthm}
\newtheorem{theorem}{Theorem}
\begin{document}
\section*{Restricted promise-capacity result}
Status: CANDIDATE pending review. P versus NP remains OPEN.
\begin{theorem}
Let a fixed public map from $S$ seed bits to $b$ code blocks use at most $t$ adaptive seed-bit queries per output coordinate. Block $i$ has length $Q_i$, code distance $d_i$, and permitted error $e_i<d_i/2$ on admitted seeds. There is one fixed heavy-cell set $V$ with
\[
 |V|\le\min\{S,(2^t-1)\sum_iQ_i/(d_i-2e_i)\}.
\]
If $A$ is the admitted set, $B=2^S-|A|$, and $R$ counts distinct complete decoded codeword tuples over $A$, then
\[
 \boxed{R\le 2^{|V|}+B.}
\]
\end{theorem}
\noindent Proof. Each output decision tree mentions at most $2^t-1$ cells, including all adaptive branches. Excluding coordinates with influence-count at least $d_i-2e_i$ makes each one-bit admitted change preserve the complete decoded codeword. Inside a fiber fixing $V$, different labels cannot be adjacent. Every nonfull label class has at least $m=S-|V|$ boundary edges in the $m$-cube, all entering bad states. Each bad state has at most $m$ incident edges. Thus a full fiber has one label and a nonfull fiber has at most as many labels as bad states. Sum over heavy-coordinate fibers.

This is an invalid-seed-count theorem, not an admission-time lower bound. A one-probe identity encoder with even-parity admissibility nearly saturates it while allowing polynomial-time promise recognition. No arbitrary-machine hardness transfer follows.
\end{document}
% END_OF_VERIFIED_REPORT::RESOURCE-DIRECTSUM::RESOURCE-PROMISE-ENTROPY12::P_VS_NP_OPEN